\documentclass[11pt]{article}
\usepackage[margin=1in]{geometry}
\usepackage[T1]{fontenc}
\usepackage{lmodern,microtype}
\usepackage{amsmath,amssymb,amsthm}
\usepackage[round,authoryear]{natbib}
\usepackage[hidelinks]{hyperref}
\newtheorem{lemma}{Lemma}
\newtheorem{theorem}{Theorem}
\title{Quality Differentiation and Market Exclusion:\\A Note on Economides (1989)}
\author{Ryota Matsuki\\Independent Researcher}
\date{Provisional manuscript --- 17 November 2026\\[2pt]\small Not approved for submission}

\begin{document}
\maketitle
\begin{abstract}
We reassess the symmetric-quality candidate behind
Economides (1989, Theorem~2). With locations held at
opposite endpoints in a covered linear Hotelling
duopoly, a finite quality increase can change
downstream price competition from sharing to
exclusion. For $1/9<\lambda<2/9$ the reduced
quality game has exactly two asymmetric exclusion
equilibria. The symmetric equilibrium also exists
if and only if $\lambda\geq4/27$. Below this
cutoff, a profitable quality deviation refutes
the proposed symmetric continuation path, but
does not imply nonexistence of pure quality
equilibrium. The correspondence is restricted to
fixed locations and does not characterize full
three-stage subgame-perfect equilibria. We do
not claim historical priority over earlier
quality-stability critiques.
\end{abstract}
\noindent\textbf{Keywords:} spatial competition; quality differentiation; Hotelling;
market exclusion; equilibrium multiplicity

\medskip
\noindent\textbf{JEL:} C72; D43; L13; R32

\section{Introduction}\label{sec:intro}
Economides \citet[pp.~27--28, eqs.~(8)--(9) and Theorem~2]{Economides1989}
reports a unique outcome combining maximal horizontal
differentiation with symmetric quality in a sequential
location--quality--price game. The calculation leading to the
symmetric quality candidate uses the local optimality
conditions for the market-sharing price regime. Such
conditions are necessary but need not suffice when a
finite quality deviation changes the subsequent price
equilibrium to market exclusion.

The logical target is the \emph{proposed equilibrium path},
not a newly solved location game. If a three-stage
equilibrium has endpoint locations and symmetric quality,
then, after those locations are chosen, each firm's
quality must be a global best response when subsequent
prices are re-optimized. A profitable quality deviation
at those same locations therefore invalidates that
proposed path. It does \emph{not} establish a replacement
subgame-perfect equilibrium of the three-stage game.

We solve the narrower question exactly. At fixed
locations $(0,1)$, for $1/9<\lambda<2/9$ the reduced
quality game has two asymmetric exclusion equilibria
throughout; the symmetric profile is a third equilibrium
if and only if $\lambda\geq4/27$. Below this cutoff, a
firm at the symmetric candidate profits by increasing
quality enough to exclude its rival. The conclusion is
not that pure quality equilibrium disappears, but that
its \emph{composition} differs from the proposed
symmetric continuation.

Earlier scholarship already identified difficulties
with the stability of quality choices in this model
\citep{Bunte1995,Bunte1997,Sorenson1995}. We provide a
self-contained, model-specific correspondence and make
no claim to have been the first to identify this issue.
Whether the exact correspondence is already contained
in \citet{Bunte1995} or \citet{Sorenson1995} remains
unresolved; this manuscript must not be submitted
without a defensible editorial assessment of that overlap.

\section{Model and equilibrium object}\label{sec:model}
There are two firms at fixed locations $\ell_1=0$ and
$\ell_2=1$ on the unit interval. Consumers $z\in[0,1]$
have uniform density $\mu>0$. The market is fully
covered: each consumer purchases from one firm,
and there is no outside option. Utility from firm $i$
at nonnegative price $p_i$ and nonnegative quality
$a_i$ is
\begin{equation}\label{eq:utility}
u_i(z)=k+a_i-p_i-|z-\ell_i|.
\end{equation}
Quality is not bounded above; zero marginal production
cost and a quadratic sunk quality cost
$ca_i^2/2$ are maintained, where $c>0$.
Write $\lambda=c/\mu$, with the parameter domain
\begin{equation}\label{eq:domain}
\frac19<\lambda<\frac29.
\end{equation}
Neither endpoint is included in our theorem.

With these locations fixed, firms first choose qualities
simultaneously, observe both qualities, and then choose
prices simultaneously. We seek pure Nash equilibria of
the reduced quality game after substituting a global
pure price equilibrium for every quality pair.
We do not optimize locations or analyze mixed strategies.

Let $\delta=a_1-a_2$. The clipped demand shares are
\begin{equation}\label{eq:demand}
s_1(p_1,p_2;\delta)=
\max\left\{0,\min\left\{1,
\frac{1+\delta-p_1+p_2}{2}\right\}\right\},
\qquad s_2=1-s_1.
\end{equation}
Dividing profit by $\mu$, firm $i$ earns
$p_i s_i-\lambda a_i^2/2$. Retaining the clipping
in \eqref{eq:demand} matters: an interior indifferent-
consumer formula cannot be used after the quality
difference makes a firm's share zero or one.

\section{Global price continuation}\label{sec:price}
\begin{lemma}[Global pure price continuation]\label{lem:price}
At fixed endpoint locations, the price equilibrium for
every pair of nonnegative qualities is
\begin{equation}\label{eq:prices}
(p_1^*,p_2^*)=
\begin{cases}
(0,-\delta-1),&\delta\leq-3,\\
(1+\delta/3,\ 1-\delta/3),&-3\leq\delta\leq3,\\
(\delta-1,0),&\delta\geq3.
\end{cases}
\end{equation}
The formulas coincide at $\delta=\pm3$.
\end{lemma}
When both firms serve positive shares, price
first-order conditions give the middle branch.
At this candidate, a full-market capture deviation
is not profitable: relative to the feasible
capture-price boundary, the normalized revenue
differences are $(\delta-3)^2/18$ for firm 1 and
$(\delta+3)^2/18$ for firm 2. For $|\delta|>3$,
the firm with higher quality instead sets the
highest price consistent with excluding its rival.
Appendix~\ref{app:priceproof} supplies the
global branch checks.

After substitution, normalized quality-stage profit
for own quality $a$ against rival quality $b$ is
\begin{equation}\label{eq:payoff}
v_\lambda(a,b)=
\begin{cases}
-\dfrac{\lambda a^2}{2},&a-b\leq-3,\\[4pt]
\dfrac12\left(1+\dfrac{a-b}{3}\right)^2
-\dfrac{\lambda a^2}{2},&-3\leq a-b\leq3,\\[6pt]
a-b-1-\dfrac{\lambda a^2}{2},&a-b\geq3.
\end{cases}
\end{equation}
All three branches have matching values at
their boundaries. Define the \emph{global}
quality best-response set as
$BR_\lambda(b)=\arg\max_{a\geq0}v_\lambda(a,b)$.

\section{Complete quality-equilibrium classification}
\label{sec:results}
Let $A=1/(3\lambda)$ and $B=1/\lambda$.
The first is the stationary symmetric-quality
candidate; the latter is the profit-maximizing
quality for a firm that captures the whole market.

\begin{theorem}[Pure quality-equilibrium correspondence]
\label{thm:complete}
For each $\lambda\in(1/9,2/9)$, the complete set
of pure Nash equilibria in the reduced quality game is
\begin{equation}\label{eq:NE}
NE(\lambda)=
\begin{cases}
\{(B,0),(0,B)\},&
  \dfrac19<\lambda<\dfrac4{27},\\[5pt]
\{(B,0),(0,B),(A,A)\},&
  \dfrac4{27}\leq\lambda<\dfrac29.
\end{cases}
\end{equation}
In particular, there are exactly two
asymmetric exclusion equilibria throughout
the interval, and an additional symmetric
equilibrium precisely when $\lambda\geq4/27$.
There are no other pure quality equilibria.
\end{theorem}

\begin{proof}[Proof outline]
Direct maximization across all three pieces
of \eqref{eq:payoff} gives
\begin{equation}\label{eq:cornerBR}
BR_\lambda(0)=\{B\},\qquad
BR_\lambda(B)=\{0\}.
\end{equation}
Thus $(B,0)$ and $(0,B)$ are equilibria,
and these exhaust equilibria containing
a zero-quality firm.

Now suppose both qualities are strictly positive
at an equilibrium. If one firm were excluded,
it would earn $-\lambda a_i^2/2<0$ and could
lower its quality to zero, remain excluded
and earn zero. This is impossible in equilibrium.
Therefore both firms sell positive shares,
their own quality choices are interior,
and the differentiable shared-market first-order
conditions require
\begin{equation}\label{eq:FOC}
(9\lambda-1)a_1+a_2=3,\qquad
a_1+(9\lambda-1)a_2=3.
\end{equation}
Subtracting gives
$(9\lambda-2)(a_1-a_2)=0$. Since
$\lambda<2/9$, any remaining equilibrium
must have $a_1=a_2=A$.

Against $A$, the sharing branch is strictly
concave and uniquely maximized at $A$,
whereas the high-quality capture branch
is uniquely maximized at the feasible $B$.
Their profit difference is
\begin{equation}\label{eq:gain}
v_\lambda(B,A)-v_\lambda(A,A)
=\frac{4-27\lambda}{18\lambda}.
\end{equation}
Consequently, $A$ is a global best response
to itself iff $\lambda\geq4/27$; equality
gives indifference and includes the symmetric
profile. Appendix~\ref{app:qualityproof}
provides the full boundary and corner checks.
\end{proof}

The economically relevant comparison is between
a locally attractive market-sharing choice and
a finite switch to exclusion. At $a-b=3$,
gross revenue is continuous, but its slope
rises from $2/3$ on the sharing side to $1$
on the exclusion side, while marginal quality
cost varies smoothly. The global deviation
gain in \eqref{eq:gain} is positive precisely
below $\lambda=4/27$. Thus lower normalized
quality cost eliminates the symmetric equilibrium
without eliminating pure equilibrium altogether;
higher cost permits the symmetric profile
without removing the exclusion equilibria.

\section{Interpretation and limitations}\label{sec:discussion}
The correction has two logically distinct parts.
First, for $\lambda<4/27$ the finite, exclusion-inducing
quality deviation is profitable against the symmetric
candidate. Because the claimed full-game path in
\citet[Theorem~2]{Economides1989} requires that
quality choice to survive optimization at its
reported endpoint locations, the proposed
symmetric path fails this necessary condition.
Second, the reduced quality game nevertheless
has two corner equilibria, $(B,0)$ and $(0,B)$;
Theorem~\ref{thm:complete} rules out additional
pure profiles by combining the two corner
best-response identities with the shared-market
first-order conditions. The new correspondence
therefore concerns quality choices \emph{conditional
on fixed endpoints}; the corner pairs are not
claims about endogenous location equilibrium.

The exclusion outcome is an equilibrium
statement about firm strategies and downstream
sales, not a welfare comparison, an empirical
prediction, or proof of a new three-stage
subgame-perfect equilibrium. Results depend on
full coverage, linear distance, quadratic
quality costs, nonnegative unbounded qualities,
nonnegative prices, and the open parameter
interval $(1/9,2/9)$. Mixed strategies,
location deviations, alternative transportation
costs and the boundary parameters are not
characterized.

The earlier general quality-stability concern
is explicitly credited to \citet{Bunte1997}.
The exact theorem-level relation to
\citet{Bunte1995} and \citet{Sorenson1995}
remains unknown because complete primary-text
comparisons are unavailable. This note neither
asserts historical priority nor treats the
absence of a priority claim as a substitute
for establishing independent publication value.

\appendix
\section{Price-stage branch verification}
\label{app:priceproof}
For $|\delta|<3$, positive-share own-price
first-order conditions yield
$(p_1,p_2)=(1+\delta/3,1-\delta/3)$;
both prices and market shares are positive.
Firm 1's candidate normalized revenue is
$\frac12(1+\delta/3)^2$. Holding $p_2$
at its candidate value, the highest
capture price is $p_2+\delta-1=2\delta/3$,
which is feasible only when nonnegative.
When feasible, the candidate-minus-capture
revenue is
\[
\frac12(1+\delta/3)^2-\frac{2\delta}{3}
=\frac{(\delta-3)^2}{18}\geq0.
\]
The symmetric comparison gives
$(\delta+3)^2/18\geq0$ for firm 2.
Other interior-price deviations cannot
beat the concave stationary revenue.
If a candidate capture price is negative,
it is not an admissible nonnegative price.
This proves global price best responses on
the central interval, including boundaries.

For $\delta>3$, take
$(p_1,p_2)=(\delta-1,0)$.
Firm 1 captures the market. Lowering
$p_1$ lowers full-market revenue, while
raising it enters the sharing region with
revenue $p_1(1+\delta-p_1)/2$,
whose unconstrained maximizer
$(1+\delta)/2$ lies below the feasible
sharing-region boundary $\delta-1$.
Thus it decreases beyond that boundary.
Firm 2, charging zero, cannot earn
positive revenue by choosing a nonnegative
price against the capture price.
A positive-share pair would require
the infeasible stationary
$p_2=1-\delta/3<0$; a priced-out firm
charging a positive price can undercut
into positive sales. Hence the stated
corner is the price equilibrium.
Interchanging the firms establishes
the branch $\delta<-3$.
At $\delta=\pm3$ the formulas
agree with the central branch.

\section{Global quality best responses}
\label{app:qualityproof}
Throughout assume $1/9<\lambda<2/9$,
so $A<3<B$.

\paragraph{The reply to zero.}
Against $b=0$, the central piece for
$0\leq a\leq3$ is
$\frac12(1+a/3)^2-\lambda a^2/2$.
Its derivative is
$1/3+(1/9-\lambda)a$,
which is decreasing but still equals
$(2-9\lambda)/3>0$ at $a=3$.
Hence that entire segment is increasing.
On $a\geq3$, the exclusive-revenue
payoff $a-1-\lambda a^2/2$
is strictly concave with unique maximum
at $B>3$. Moreover,
\[
v_\lambda(B,0)-v_\lambda(3,0)
=\frac{(1-3\lambda)^2}{2\lambda}>0.
\]
Therefore $BR_\lambda(0)=\{B\}$.

\paragraph{The reply to $B$.}
Against $b=B$, any $0<a\leq B-3$
is excluded, earns $-\lambda a^2/2<0$,
and is inferior to $a=0$ with payoff zero.
On $a\in[B-3,B+3]$, the sharing-branch
profit is strictly negative at the left
endpoint, with value
$-\lambda(B-3)^2/2$.
Its first derivative there is
$-\lambda(B-3)<0$, and the second
derivative is $1/9-\lambda<0$;
it therefore decreases throughout.
On $a\geq B+3$, the reverse-exclusion
payoff $a-B-1-\lambda a^2/2$
is decreasing, as its unconstrained
maximizer $B$ lies below the branch,
and its boundary value is already
negative. Thus every positive own
quality loses to zero and
$BR_\lambda(B)=\{0\}$.

\paragraph{The reply to $A$.}
Since $A<3$, the low-exclusion branch
against $A$ cannot occur at $a\geq0$.
The strictly concave central branch
on $[0,A+3]$ is maximized at $A$.
The high-exclusion branch on
$[A+3,\infty)$ is maximized at
$B$, which is feasible because
$B-A=2/(3\lambda)>3$.
The two candidate payoffs are
\[
v_\lambda(A,A)=\frac12-\frac{1}{18\lambda},
\qquad
v_\lambda(B,A)=\frac{1}{6\lambda}-1.
\]
The difference is exactly
\eqref{eq:gain}. For $\lambda=4/27$,
both $A$ and $B$ best respond to $A$,
but $(B,A)$ cannot be an equilibrium
because the sole best response to
$B$ is zero.

\paragraph{Exhaustive necessity.}
If both qualities are positive,
a firm with zero sales can decrease
its quality to zero without obtaining
negative sales, saving a positive
sunk cost. Thus both firms must
sell positive shares; the interior
quality FOCs \eqref{eq:FOC}
then force $a_1=a_2=A$.
All remaining cases contain
a zero-quality choice and are
fully covered by the two corner
best-response identities. This
establishes necessity and sufficiency.

\section{Formal scope and research transparency}
\label{app:transparency}
For $\lambda>0$, the normalization
$x=\lambda a$, $y=\lambda b$ is
a bijection on nonnegative qualities.
The strictly positive payoff scaling
$18\lambda^2v_\lambda(x/\lambda,y/\lambda)$
preserves all global best responses.
The Lean theorem
\texttt{ShortNoteFormal.complete\_pure\_quality\_NE}
certifies an if-and-only-if Nash
classification for that \emph{encoded
reduced quality payoff}, over all
nonnegative deviations. It does not
derive clipped consumer demand or
the three-branch price equilibrium from
primitive preferences, and does not
prove equilibrium of the full
location--quality--price game.
The code and exact-rational regression
checks are available in the linked
research repository with the
corresponding formal certificates.

The human-review and generative-AI
declaration is provided separately below.
The mathematical theorem is a working
pre-freeze candidate: current project
governance has not issued canonical
theory-freeze or submission approval.

\section*{Declaration of generative-AI assistance}
OpenAI's ChatGPT was used during 2026 to assist
theoretical exploration, algebraic and
computational checking, proof formalization,
literature discovery and organization, and
manuscript drafting. The author reviewed the
model-specific equilibrium classification and
its case-by-case proof, and remains responsible
for the mathematical claims, references,
originality assessment and final text.
The Lean formalization checks the encoded
reduced-payoff Nash correspondence; it
does not independently establish the
primitive price game or the original
three-stage equilibrium.
Complete version-level AI-use details and
the author's final publication declaration
must be reconciled with the research log
and Wiley's current policy before submission.

\bibliographystyle{apalike}
\bibliography{../../references/references}
\end{document}
